\newcommand{\dodecahedron}[3][(0,0,0)]{
  \tdplotsetmaincoords{#2}{#3}
  \begin{tikzpicture}[
    maincolor/.style=red,
    scale=1.7,
    tdplot_main_coords,
    axis/.style={-stealth},
    mainline/.style={maincolor,thick},
    invisibleline/.style={white,densely dashed},
    ]
%
    \tikzmath{
      \kphi = 1.618033988749895;
      \kihp = 1 / \kphi;
      \nphi = -1*\kphi;
      \nihp = -1*\kihp;
    }
%    axis
    % \draw[axis](-1.7,0,0)--(1.7,0,0)node[pos=1.1]{$x$};
    % \draw[axis](0,-1.7,0)--(0,1.7,0)node[pos=1.1]{$y$};
    % \draw[axis](0,0,-1)--(0,0,2)node[pos=1.1]{$z$};
%
    % \coordinate[shift={#1}] (O) at   (0,0,0);
    % \coordinate[shift={#1}] (p1) at  (0,1,0);
    % \coordinate[shift={#1}] (p2) at  (0,0,1);
    % \coordinate[shift={#1}] (p3) at  (1,1,0);
    % \coordinate[shift={#1}] (p4) at  (0,1,1);
%
    %     % vertices
    % a cube: (1,2,3,4), (5,6,7,8), (1,2,5,6), (3,4,7,8), (1,3,5,7), (2,4,6,8)
    \coordinate[shift={#1}] (q1) at  (1,   1,  1);
    \coordinate[shift={#1}] (q2) at  (1,   1, -1);
    \coordinate[shift={#1}] (q3) at  (1,  -1,  1);
    \coordinate[shift={#1}] (q4) at  (1,  -1, -1);
    \coordinate[shift={#1}] (q5) at  (-1,  1,  1);
    \coordinate[shift={#1}] (q6) at  (-1,  1, -1);
    \coordinate[shift={#1}] (q7) at  (-1, -1,  1);
    \coordinate[shift={#1}] (q8) at  (-1, -1, -1);
    % vertices that overlap with a cube:
    %
    % (1,2), (2,3), (3,4), (4,1)
    % (1,5), (2,6), (3,7), (4,8)
    % (4,5), (5,6), (6,7), (7,8)
    %
    % then find the vertices closest to each pair of cube vertices
    %
    % * a pentagon point will be at (0,phi,phi^-1).
    % * it's tail at (+phi^-1,0,-phi^-1)
    %
    \coordinate[shift={#1}] (p1) at  (0,      \kphi, \kihp);
    \coordinate[shift={#1}] (p2) at  (0,      \nphi, \kihp);
    \coordinate[shift={#1}] (p3) at  (0,      \nphi, \nihp);
    \coordinate[shift={#1}] (p4) at  (0,      \kphi, \nihp);
%   
    \coordinate[shift={#1}] (p5) at  (\kihp,  0,     \kphi);
    \coordinate[shift={#1}] (p6) at  (\nihp,  0,     \kphi);
    \coordinate[shift={#1}] (p7) at  (\nihp,  0,     \nphi);
    \coordinate[shift={#1}] (p8) at  (\kihp,  0,     \nphi);
 %  
    \coordinate[shift={#1}] (p9) at  (\kphi,  \kihp, 0);
    \coordinate[shift={#1}] (p10) at (\nphi,  \kihp, 0);
    \coordinate[shift={#1}] (p11) at (\nphi,  \nihp, 0);
    \coordinate[shift={#1}] (p12) at (\kphi,  \nihp, 0);
% %
%     % fill faces
    \fill[maincolor,opacity=0.1](p1)--(p2)--(p3)--(p4)--(p9)--cycle;
    \fill[maincolor,opacity=0.1](p5)--(p6)--(p7)--(p8)--(p9)--cycle;
    % \fill[maincolor,opacity=0.1](p1)--(p2)--(p3)--(p4)--(p12)--cycle;
    % \fill[maincolor,opacity=0.1](p2)--(p3)--(p4)--cycle;
    % \fill[maincolor,opacity=0.1](p3)--(p4)--(p5)--cycle;
    % \fill[maincolor,opacity=0.1](p1)--(p2)--(p3)--cycle;
    % \fill[maincolor,opacity=0.1](p2)--(p3)--(p4)--cycle;
    % \fill[maincolor,opacity=0.1](p3)--(p4)--(p5)--cycle;
    
    % \fill[maincolor,opacity=0.1](p4)--(p5)--(p6)--cycle;
%
    % % inner product
    % \pgfmathparse{-sin(\tdplotmaintheta)*sin(\tdplotmainphi)+(1+cos(60))*sin(\tdplotmaintheta)*cos(\tdplotmainphi)/sin(60)+cos(\tdplotmaintheta)/sqrt(2)>0}
    % \pgfmathsetmacro{\OBCprod}{\pgfmathresult}
    % \pgfmathparse{-sin(\tdplotmaintheta)*sin(\tdplotmainphi)-(1+cos(60))*sin(\tdplotmaintheta)*cos(\tdplotmainphi)/sin(60)+cos(\tdplotmaintheta)/sqrt(2)>0}
    % \pgfmathsetmacro{\OABprod}{\pgfmathresult}
    % \pgfmathparse{sin(\tdplotmaintheta)*sin(\tdplotmainphi)/cos(60)+cos(\tdplotmaintheta)/sqrt(2)>0}
    % \pgfmathsetmacro{\OCAprod}{\pgfmathresult}
    % \pgfmathparse{-cos(\tdplotmaintheta)>0}
    % \pgfmathsetmacro{\ABCprod}{\pgfmathresult}

    % \ifthenelse{\ABCprod=0 \and \OABprod=0 \and \OBCprod=0 \and \OCAprod=1}{
    % \draw[invisibleline](O)--(B) (A)--(B) (C)--(B); \draw[mainline](A)--(O)--(C)--cycle; % 0001
    % }{
    %   \ifthenelse{\ABCprod=0 \and \OABprod=0 \and \OBCprod=1 \and \OCAprod=0}{
    %   \draw[invisibleline](O)--(A) (B)--(A) (C)--(A); \draw[mainline](O)--(B)--(C)--cycle; % 0010
    % }{
    %   \ifthenelse{\ABCprod=0 \and \OABprod=0 \and \OBCprod=1 \and \OCAprod=1}{
    %   \draw[mainline](O)--(C) (O)--(A)--(C) (O)--(B)--(C);\draw[invisibleline](A)--(B); % 0011
    % }{
    %   \ifthenelse{\ABCprod=0 \and \OABprod=1 \and \OBCprod=0 \and \OCAprod=0}{
    %   \draw[invisibleline](O)--(C) (A)--(C) (C)--(B); \draw[mainline](A)--(B)--(O)--cycle; % 0100
    % }{
    %   \ifthenelse{\ABCprod=0 \and \OABprod=1 \and \OBCprod=0 \and \OCAprod=1}{
    %   \draw[mainline](O)--(A) (O)--(B)--(A) (O)--(C)--(A);\draw[invisibleline](B)--(C); % 0101
    % }{
    %   \ifthenelse{\ABCprod=0 \and \OABprod=1 \and \OBCprod=1 \and \OCAprod=0}{
    %   \draw[mainline](O)--(B) (O)--(A)--(B) (O)--(C)--(B);\draw[invisibleline](A)--(C); % 0110
    % }{
    %   \ifthenelse{\ABCprod=0 \and \OABprod=1 \and \OBCprod=1 \and \OCAprod=1}{
    %   \draw[mainline](O)--(A) (O)--(B) (O)--(C) (A)--(B)--(C)--cycle; % 0111
    % }{
    %   \ifthenelse{\ABCprod=1 \and \OABprod=0 \and \OBCprod=0 \and \OCAprod=0}{
    %   \draw[invisibleline](O)--(A) (O)--(B) (O)--(C); \draw[mainline](A)--(B)--(C)--cycle; % 1000
    % }{
    %   \ifthenelse{\ABCprod=1 \and \OABprod=0 \and \OBCprod=0 \and \OCAprod=1}{
    %   \draw[mainline](A)--(C) (A)--(O)--(C) (A)--(B)--(C);\draw[invisibleline](O)--(B); % 1001
    % }{
    %   \ifthenelse{\ABCprod=1 \and \OABprod=0 \and \OBCprod=1 \and \OCAprod=0}{
    %   \draw[mainline](B)--(C) (B)--(A)--(C) (B)--(O)--(C);\draw[invisibleline](A)--(O); % 1010
    % }{
    %   \ifthenelse{\ABCprod=1 \and \OABprod=0 \and \OBCprod=1 \and \OCAprod=1}{
    %   \draw[mainline](O)--(A) (O)--(B) (O)--(C) (A)--(B)--(C)--cycle; % 1011
    % }{
    %   \ifthenelse{\ABCprod=1 \and \OABprod=1 \and \OBCprod=0 \and \OCAprod=0}{
    %   \draw[mainline](A)--(B) (A)--(C)--(B) (A)--(O)--(B);\draw[invisibleline](C)--(O); % 1100
    % }{
    %   \ifthenelse{\ABCprod=1 \and \OABprod=1 \and \OBCprod=0 \and \OCAprod=1}{
    %   \draw[mainline](O)--(A) (O)--(B) (O)--(C) (A)--(B)--(C)--cycle; % 1101
    % }{
    %   \ifthenelse{\ABCprod=1 \and \OABprod=1 \and \OBCprod=1 \and \OCAprod=0}{
    %   \draw[mainline](O)--(A) (O)--(B) (O)--(C) (A)--(B)--(C)--cycle; % 1110
    % }{%
    % }}}}}}}}}}}}}}
%
%   draw points
  % \foreach \P in{p1,p2,p3,p4,p5,p6,p7,p8,p9,p10,p11,p12}
  % \path(0,0,{sqrt(2)/3})--(\P)node[pos=1.2]{\P};
  \end{tikzpicture}
  % }
}
